\programchapter{Code of Conduct [EN]}

The organizers are committed to making this meeting productive and enjoyable for
everyone, regardless of gender, sexual orientation, disability, physical appearance,
body size, race, nationality, religion, career stage, or institutional affiliation.
Harassment, discrimination, intimidation, bullying, or retaliation will not be
tolerated at any conference event.

Please follow these guidelines:

\begin{itemize}
  \item \textbf{Behave professionally:} sexist, racist, exclusionary, or otherwise
  demeaning comments or jokes are not appropriate. Harassment includes sustained
  disruption of talks or events, inappropriate physical contact, sexual attention
  or innuendo, deliberate intimidation, stalking, and photography or recording of
  an individual without consent.
  \item \textbf{Communicate appropriately:} all discussion should be suitable for a
  professional audience that includes people from many different backgrounds.
  Sexual language or imagery is not appropriate in talks, posters, discussions,
  social media connected to the meeting, or informal conference spaces.
  \item \textbf{Treat others with respect:} be constructive, give proper credit,
  and do not insult, belittle, or deliberately exclude other participants.
\end{itemize}

Participants who are asked to stop inappropriate behaviour are expected to comply
immediately. Organizers reserve the right to take any action they deem necessary,
including removal from the meeting without refund.

Any participant who wishes to report a concern is encouraged to speak in
confidence with one of the following:

\begin{itemize}
  \item an LOC chair or co-chair
  \item a member of the CASCA Equity, Diversity and Inclusion Committee
  \item a CASCA Board member or ombudsperson
\end{itemize}

\textbf{Protocol}

\begin{itemize}
  \item The designated CASCA contact will request a written record of the concern,
  including the time, date, and relevant particulars, and will confirm the record
  with the reporting participant.
  \item The concern will be brought to the attention of the LOC and CASCA
  leadership so that an appropriate response can be determined promptly.
  \item The person or people named in the concern may be informed of the allegation
  and invited to provide their account of events.
\end{itemize}

This meeting guidance follows the structure of the CASCA 2025 annual meeting
policy and is intended to be read alongside the official
\href{https://casca.ca/?page_id=22442}{CASCA Code of Conduct}.